
\section{Introducció}

\lettrine[lines=3]{E}{n} els darrers anys, l'avenç de la Intel·ligència Artificial (IA) ha estat notori, transformant radicalment la nostra capacitat per processar informació i automatitzar tasques complexes. Aquesta eclosió tecnològica ofereix una oportunitat sense precedents per millorar diversos aspectes de la societat, aportant eficiència i precisió en àmbits que fins ara depenien de processos manuals i lents.

Dins d'aquest paradigma, es proposa el present projecte com una solució innovadora per al manteniment del Campus. L'objectiu és garantir que les instal·lacions es mantinguin en bon estat mitjançant una plataforma que aprofiti la potència de la IA per a la classificació d'incidències i l'anàlisi de prioritats basant-se en imatges. Aquesta aplicació no només busca l'eficiència operativa, sinó que vol millorar substancialment l'experiència dels alumnes, convertint-los en agents actius del seu entorn.

L'adopció d'una solució àgil i flexible com la que es presenta comporta avantatges significatius: la capacitat de resposta en temps real , la transparència en el seguiment de les incidències mitjançant una màquina d'estats definida i la facilitat de report des de dispositius mòbils amb geolocalització. La implementació d'aquest sistema suposarà un gran canvi en la gestió del Campus, modernitzant la comunicació entre la comunitat universitària i els equips de manteniment.

\section{Estat de l'art}
La gestió d’infraestructures en grans complexos com un campus universitari ha estat tradicionalment un procés reactiu i unidireccional. Històricament, la detecció i resolució de desperfectes s'ha basat en formularis web estàtics, correus electrònics a centres de suport o trucades telefòniques. Aquest model genera colls d'ampolla en l'administració degut al triatge manual, i un efecte de "caixa negra" per a l'usuari, el qual rarament rep retroalimentació sobre l'estat de la seva petició. El present projecte s’emmarca en la necessitat de modernitzar aquest flux, analitzant les solucions actuals del mercat i determinant quins aspectes diferencials s'aporten per transformar la UAB en un entorn col·laboratiu i intel·ligent.

\subsection{Evolució dels sistemes tradicionals}
Fins fa poc, el programari de manteniment es limitava a grans sistemes de gestió de manteniment limitats, dissenyades exclusivament per a tècnics i administradors. Aquests sistemes han evolucionat a sistemes mes flexibles i usables per als usuaris aprofitant l'us dels telefons movils i sistemes de comunicacions moderns.

Avui dia, els dispositius mòbils no només permeten la comunicació, sinó que actuen com a sensors distribuïts gràcies al maquinari integrat (càmeres d'alta resolució, mòduls GPS de gran precisió). Les empresen han aprofitat aquesta onada tecnològica per poder desenvolupar aplicacions que permeten gestionar i processar les incidencies de les infraestructures de manera remota. Aquests sistemes permeten als usuaris crear incidències amb comoditat i veure el seguiment d'aquestes. Aquest projecte presenta una arquitectura híbrida (Web per a administració, App per a usuaris/tècnics) que descentralitza la detecció d'incidències i n'automatitza el processament, reduint dràsticament els temps de resposta.

\subsection{Anàlisi de Solucions Existents}
Per situar el projecte en el panorama tecnològic actual, s’ha realitzat un anàlisi exhaustiu de dos models de referència: l'aplicació "Find It, Fix It" de la ciutat de Seattle i el model de mobilitat compartida de Bolt CityFix. Aquesta comparativa permet identificar els buits operatius que la present plataforma pretén resoldre.

\begin{itemize}
	\item \textbf{Find It, Fix It (Ciutat de Seattle)}\cite{10}: Aquest sistema actua sota una arquitectura d’integració en sistemes empresarials complexos, permet gestionar les incidencies de la ciutat de Seattle. Permet adjuntar evidències visuals a les incidencies i utilitza metadades d'ubicació per centralitzar els reports. Permet visualitzar graficament en un mapa on s'ubiquen les incidencies. 
	Tot i aixi, no hi ha cap plataforma externa web d'administració i els usuaris sovint perceben una "manca d'acció" quan els seus casos es tanquen prematurament a causa de redireccionaments interns entre departaments.
	
	\item \textbf{Bolt CityFix}\cite{11}: Bolt és una superaplicació europea de mobilitat que ofereix serveis de transport sota demanda (VTC i taxi), lloguer de patinets/bicicletes elèctriques i repartiment de menjar. Bolt aprofita la seva xarxa de micromobilitat preexistent per convertir els conductors i usuaris en sensors actius de la ciutat mitjançant el report d'incidencies i millores.
	
	Destaca per la immediatesa i el volum de reports (forats, senyals, semàfors) inabastable per a un ciutadà mitjà, integrant l'eina dins d'una "super-app" de mobilitat. 
	El seu enfocament és eminentment reactiu i centrat en la infraestructura de transport, per aixó manca d'una interfície d'administració profunda per a tècnics.
\end{itemize}

\subsection{Proposta de valor}
A diferència de les solucions estàndard de mercat, aquest projecte es posiciona com una plataforma activa i intel·ligent, reforçant en tres aspectes que resolen les mancances de les aplicacions de gestió actuals:

\begin{itemize}
	\item \textbf{Gestió d'incidencies amb agents i IA}: Mitjançant la integració de models de llenguatge (LLMs) via LangChain i OpenAI, el sistema no és un simple gestor de base de dades. Els agents d'IA analitzen el llenguatge natural de l'usuari, n'avaluen la criticitat, filtren duplicats mitjançant anàlisi de similitud i generen notificacions automatitzades. Això descarrega el personal administratiu de tasques repetitives. Aquest aspecte té molt de potencial ja que amb l'implementació de agents es pot escalar l'aplicació automatitzant la majoria de funcionalitats. Fins arribar al punt de crear una plataforma d'incidencies completament automatica que actui sota la revisió i control humà.   
	
	\item \textbf{Interficie web d'administració}: Aquest projecte implementa una plataforma tant movil com d'escitori, permetent als usuaris administradors la visualització de dades i grafics en una interficie comoda d'escritori. Gràcies a l'ús de WebSockets, les mètriques dels garfics i mapes s'actualitzen instantàniament, permetent una coordinació d'emergències altament eficient.
	
	\item \textbf{Sistema de recompenses (Gamificació)}: Per combatre el report deficient o malintencionat, s'implementa un mecanisme de puntuació on les bones pràctiques es premien (fomentant l'esperit de pertinença i cura del campus), mentre que el mal ús deliberat està vinculat a sancions acadèmiques, garantint així l'alta fiabilitat de les dades d'entrada.
	
\end{itemize}

\section{Objectius}
El projecte present es planteja amb l’objectiu principal de resoldre la necessitat de manteniment i millora de les infraestructures del campus universitari. Mitjançant tecnologies de la informació serem capaços de notificar als estudiants, mantenir i gestionar el Campus en bones condicions per garantir una bona experiència als alumnes i professors. Els objectius es centren a maximitzar l'eficiència del manteniment i potenciar el sentiment de comunitat.


\begin{itemize}
	\item \textbf{Optimització de la disponibilitat dels recursos del campus}: Reduir dràsticament el temps de permanència de les infraestructures en mal estat (com ara ordinadors de sales d'estudi, mobiliari d'aules o il·luminació). L'objectiu és que qualsevol element defectuós sigui detectat i reportat en qüestió de minuts, minimitzant l'impacte en l'activitat acadèmica.
	
	\item \textbf{Foment de la co-responsabilitat i la participació ciutadana}: Transformar l'alumne de "subjecte passiu" a "agent actiu" del manteniment del campus. Mitjançant el sistema de recompenses i gamificació, s'estimula un comportament cívic i proactiu, creant una comunitat que vetlla pel seu entorn.
	
	\item \textbf{Millora de l'eficiència operativa dels equips tècnics}: Proporcionar als operaris de manteniment una eina que prioritzi les tasques segons la gravetat i la ubicació. Això permet una assignació de recursos més intel·ligent, passant d'un model de manteniment reactiu a un model més organitzat i basat en dades reals.
	
	\item \textbf{Anàlisi de dades per a la gestió estratègica del Campus}: Utilitzar la informació geolocalitzada recollida per identificar problemes estructurals o zones amb una degradació superior a la mitjana. Això servirà de suport a la presa de decisions de la universitat i al control de l'estat actual del Campus.
	
	\item \textbf{Reducció de costos administratius mitjançant l'automatització}: Delegar en la Intel·ligència Artificial el triatge i la classificació d'incidències per reduir la càrrega de treball manual dels gestors. L'objectiu és accelerar la burocràcia interna i permetre que els recursos humans es concentrin en les tasques essencials.
	
\end{itemize}

\section{Metodología}
Per a la realització d'aquest projecte s'ha optat per una metodologia de desenvolupament àgil basada en el marc de treball Scrum\cite{4}, adaptada a la naturalesa individual d'un Treball de Final de Grau. Aquest enfocament permet una gestió iterativa i incremental, assegurant que les funcionalitats més crítiques es trobin en una base sòlida abans d'afegir capes de complexitat.

Per dur a terme aquesta metodología, en primer lloc es defineix el \textit{Product Backlog} a partir dels requisits generats i funcionalitats definides. A partir d'aquí, s'extreu el \textit{Sprint Backlog}, que representa el subconjunt de tasques seleccionades per a ser executades en cada \textit{Sprint}. Finalment, cada cicle culmina amb un Increment, que és el resultat tangible i funcional del programari que s'afegeix de forma acumulativa a les versions anteriors, garantint així una evolució del projecte constant, robusta i verificable en cada lliurament.

En aquest cas, en lloc de fer reunions diàries es farà un seguiment setmanal amb el tutor del TFG (\textit{Sprint Review i Retrospective}) per avaluar l'increment generat.
Aquesta metodologia és ideal per a aquest projecte perquè permet integrar components complexos (com la IA i els SIG) de manera progressiva, garantint que la base del sistema sigui sòlida abans d'afegir-hi capes addicionals.


\section{Planificació}
Per establir una referència temporal clara i assolible, la planificació es divideix en blocs de treball de dues setmanes de durada, anomenats Sprints. Aquest enfocament permet fragmentar la complexitat del sistema en etapes consecutives, on cada cicle té com a objectiu el lliurament d'un increment funcional del programari. En aquest projecte hem definit Sprints de dues setmanes per tenir mes marge de desenvolupament. A continuació es mostra el diagrama de Gantt complet del projecte:

% Per a fer que la figura ocupi les dues columnes utilitzeu "figure*" per comptes de "figure"
\begin{figure}[!h]
	\centering
	\includegraphics[width=0.4\textwidth]{img/gantt}
	\caption{Diagrama de Gantt}
	\label{fig-exemple}
\end{figure}

Els Sprints que es durant a terme en aquest projecte son el següents:

\begin{flushleft}
	\textbf{Sprint 1: Definició i Disseny d'Arquitectura. }\\
\end{flushleft}
Aquest primer cicle es focalitza en la concreció dels objectius, l'anàlisi detallada dels requisits i el disseny de l'esquema de dades relacional. Es defineix l'arquitectura del sistema i el model conceptual de la màquina d'estats que governarà les incidències.

\begin{flushleft}
	\textbf{Sprint 2: Infraestructura Backend i Lògica de Control.} \\
\end{flushleft}
Durant aquesta fase es realitza la configuració del servidor amb Node.js i l'ORM Prisma. S'implementa la base de dades PostgreSQL amb l'extensió PostGIS per a la gestió geogràfica i es programa la lògica de transicions d'estats amb XState.

\begin{flushleft}
	\textbf{Sprint 3: Desenvolupament de la Interfície Web i SIG. } \\
\end{flushleft}
Es construeix el dashboard d'administració amb React, integrant les APIs de cartografia per a la visualització de clústers i mapes de calor. S'estableixen les eines de filtratge i assignació manual de tasques.

\begin{flushleft}
	\textbf{Sprint 4: Mòdul de l'Alumne i Aplicació Mòbil. } \\
\end{flushleft}
Inici del desenvolupament amb React Native i Expo\cite{3}. Es crea el formulari de report d'incidències, integrant l'accés natiu a la càmera i la captura de coordenades GPS.

\begin{flushleft}
	\textbf{Sprint 5: Mòdul Tècnic i Comunicació en Temps Real. } \\
\end{flushleft}
Es desenvolupa la interfície específica per a operaris amb llistats per proximitat. S'implementen les notificacions push i la connexió bidireccional amb Socket.io per a l'actualització instantània del sistema.

\begin{flushleft}
	\textbf{Sprint 6: Integració d'Intel·ligència Artificial.} \\
\end{flushleft}
En aquesta etapa s'integren OpenAI i LangChain al backend per permetre la classificació automàtica d'incidències mitjançant imatges i la generació intel·ligent d'avisos crítics.

\begin{flushleft}
	\textbf{Sprint 7: Gamificació i Control de Qualitat (QA).} \\
\end{flushleft}
Es programa la lògica del sistema de recompenses i rànquings. Es realitza una fase intensiva de proves de sistema i correcció d'errors per garantir la robustesa de la plataforma i el compliment del RGPD.

\begin{flushleft}
	\textbf{Sprint 8: Preparació de la defensa} \\
\end{flushleft}
El darrer cicle es dedica al tancament de la memòria escrita, el refinament final de l'experiència d'usuari (UI/UX) i la preparació dels materials per a la defensa pública del projecte.


\section{Requisits}
Per garantir l'èxit en el disseny i desenvolupament de la plataforma integral de gestió de desperfectes del campus universitari, s'han identificat i classificat els següents requisits funcionals i no funcionals. Aquests requisits delimiten l'abast tècnic del projecte i satisfan les necessitats dels diferents actors involucrats.

\subsection{Requisits Funcionals}
Aquests requisits defineixen les capacitats d'interacció i el comportament del sistema davant les accions dels usuaris:

\begin{itemize}
	\item \textbf{RF-01: Gestió de rols i privilegis. }
	El sistema ha de permetre l'accés a tres perfils diferenciats: Administrador (visió global i control), Tècnic (resolució operativa) i Alumne (report d'incidències).
	
	\item \textbf{RF-02: Autenticació d'usuaris.} 
	El sistema ha de permetre en el registre i inici de sessió la integració amb sistemes d'autenticació de Google i Outlook per a aquells que disposin de compte de l'autonoma. 
	
	\item \textbf{RF-03: Creació d'incidències mòbils. }
	Els alumnes han de poder registrar incidències des de l'app amb estat "Oberta", adjuntant descripció, coordenades GPS i proves fotogràfiques.
	
	\item \textbf{RF-04: Control del flux d'estats (XState).} 
	El sistema ha de governar el cicle de vida de la incidència mitjançant 5 estats finits: Oberta, Assignada, En procés, Validada i Tancada 
	
	\item \textbf{RF-05: Restricció de transicions.} 
	Només es pot assignar un usuari amb rol "Tècnic" per moure una incidència que esta "Oberta" a l'estat "Asignada”. 
	
	\item \textbf{RF-06: Traçabilitat d'actors.} 
	Cada incidència ha d'estar vinculada obligatòriament a un ID d'usuari (creador) i a un ID de tècnic (assignat). 
	
	\item \textbf{RF-07: Panell de control d'administració.} 
	Els administradors han de poder visualitzar llistats, filtrar incidències i assignar-les a tècnics recomanats.
	
	\item \textbf{RF-08: Gestió operativa per a tècnics.} 
	Els tècnics han de poder autoassignar-se tasques mitjançant un llistat filtrat per proximitat geogràfica. 
	
	\item \textbf{RF-09: Verificació fotogràfica de resolució.} 
	Els tècnics han de reportar la solució adjuntant fotos de l'abans i el després per canviar l'estat a "Validada”. 
	
	\item \textbf{RF-10: Tancament d'incidències.} 
	L'administrador ha de rebre les peticions  "Validades" i determinar si es tanquen definitivament o es reassignen. 
	
	\item \textbf{RF-11: Visualització cartogràfica (SIG).} 
	El sistema ha de representar les incidències en un mapa interactiu amb clustering i mapes de calor per zones.
	
	\item \textbf{RF-12: Notificacions Push. }
	Enviament d'alertes als dispositius mòbils d'usuaris i tècnics davant canvis d'estat o avisos crítics. 
	
	\item \textbf{RF-13: Actualització en temps real.} 
	El tauler d'administració s'ha d'actualitzar instantàniament en rebre nous reports o canvis de disponibilitat de tècnics. 
	
	\item \textbf{RF-14: Classificació automàtica per IA. }
	Ús d'agents de IA per analitzar imatges, classificar el tipus d'incidència i determinar la prioritat. 
	
	\item \textbf{RF-15: Redacció d'avisos per IA. }
	Agent especialitzat per redactar notificacions personalitzades i alertes en cas d'incidències crítiques. 
	
	\item \textbf{RF-16: Sistema de Recompenses i Sancions. }
	Implementació d'un rànquing de punts per a usuaris col·laboradors i aplicació de sancions acadèmiques per mal ús.
	
\end{itemize}

\subsection{Requisits No Funcionals}

\begin{itemize}
	
	\item \textbf{RNF-01: Usabilitat i Experiència d'Usuari (UX/UI).} El sistema ha de presentar una interfície intuïtiva i adaptada a cada plataforma (mòbil per a l'ús en moviment i web per a l'anàlisi de dades), garantint que un alumne pugui completar un report d'incidència en menys de 30 segons.
	
	\item \textbf{RNF-02: Disponibilitat Multiplataforma.} La lògica de negoci ha d'estar centralitzada al backend per garantir la consistència de les dades, sent accessible tant des de navegadors web moderns com des de dispositius iOS i Android (vía Expo).
	
	\item \textbf{RNF-03: Rendiment i Eficiència Geoespacial.} El sistema ha de ser capaç de processar consultes espacials complexes (proximitat i clústers) mitjançant PostGIS amb un temps de resposta inferior a 500ms per a l'usuari final.
	
	\item \textbf{RNF-04: Accés a Recursos Natius.} L'aplicació mòbil ha d'integrar-se de forma transparent amb el maquinari del dispositiu (GPS i càmera), assegurant la captura de coordenades amb una precisió mínima de 10 metres.
	
	\item \textbf{RNF-05: Seguretat i Control d'Accés (RBAC).} S'ha d'implementar un control d'accés basat en rols (Role-Based Access Control) que restringeixi les operacions de la base de dades i les rutes de l'API segons el perfil de l'usuari autenticat.
	
	\item \textbf{RNF-06: Seguretat, Privacitat i Compliment (RGPD).} El sistema ha de complir el Reglament General de Protecció de Dades, incloent l'encriptació de dades sensibles en repòs i el trànsit de dades mitjançant protocols segurs (HTTPS/WSS).
	
	\item \textbf{RNF-07: Baixa Latència i Temps Real.} Les actualitzacions del panell d'administració s'han de reflectir mitjançant \textit{Socket.io} amb una latència gairebé instantània (idealment < 200ms) per garantir la coordinació operativa dels tècnics.
	
	\item \textbf{RNF-08: Mantenibilitat i Qualitat del Codi.} El sistema s'ha de desenvolupar seguint patrons de disseny modulars i utilitzant un ORM (Prisma) que faciliti la traçabilitat de les migracions i la validació de tipus (TypeScript), assegurant que el programari sigui fàcil d'ampliar en el futur.
	
\end{itemize}


\section{Eines Utilitzades}

En aquesta secció es presenten les principals eines i tecnologies emprades durant el desenvolupament del projecte CityFix, agrupades per àmbit funcional.

\subsection{Llenguatges i entorns d'execució}

\begin{itemize}
	\item \textbf{TypeScript}: Superconjunt tipat de JavaScript utilitzat tant al backend com al frontend. El tipatge estàtic permet detectar errors en temps de compilació i millora la mantenibilitat del codi en projectes de mida mitjana-gran.
	\item \textbf{Node.js}: Entorn d'execució de JavaScript al servidor. Permet compartir el mateix llenguatge entre frontend i backend, reduint el canvi de context durant el desenvolupament.
	\item \textbf{SQL / PL / pgSQL}: Utilitzat per a les consultes avançades que Prisma no suporta nativament i per als triggers de PostGIS que sincronitzen les coordenades amb la columna de geometria. 
\end{itemize}

\subsection{Backend}

\begin{itemize}
	\item \textbf{Express v5}: Framework HTTP minimalista per a Node.js. Proporciona routing, middlewares encadenables i gestió de peticions/respostes. 
	
	\item \textbf{Prisma v7}: ORM (Object-Relational Mapping) per a TypeScript que genera un client tipat a partir de l'esquema de la base de dades.\cite{7}
	
	\item \textbf{node-postgres (pg)}: Driver natiu de PostgreSQL per a Node.js. Utilitzat com a adaptador subjacent de Prisma v7 mitjançant \texttt{@prisma/adapter-pg}, proporcionant la connexió directa amb la base de dades i permetent l'execució de consultes SQL. \cite{8}
	
	\item \textbf{XState v5}: Llibreria de màquines d'estats finits i statecharts. S'utilitza per governar el cicle de vida de les incidències amb l'API \texttt{setup()}, que permet definir guards amb noms tipats per al control d'accés per rol.
	
	\item \textbf{jsonwebtoken (JWT)}: Llibreria per a la generació i verificació de tokens JSON Web Token. Permet implementar autenticació stateless amb signatura criptogràfica i expiració configurable.
	
	\item \textbf{bcrypt}: Algoritme de hashing de contrasenyes basat en Blowfish. Genera un salt aleatori per cada hash, garantint que contrasenyes idèntiques produeixin resultats diferents.
	
	\item \textbf{dotenv}: Mòdul que carrega variables d'entorn des d'un fitxer \texttt{.env} al procés de Node.js. 
	
	\item \textbf{cors}: Middleware per a Express que configura les capçaleres Cross-Origin Resource Sharing, permetent que el frontend accedeixi al backend des d'un origen diferent en producció.
	
\end{itemize}

\subsection{Frontend web}

\begin{itemize}
	\item \textbf{React 19}: Llibreria d'interfícies d'usuari basada en components declaratius. Permet construir SPAs (Single Page Applications) amb un model de dades reactiu i un cicle de vida gestionat per hooks. \cite{12}
	
	\item \textbf{Vite}: Bundler i servidor de desenvolupament amb compilació instantània i Hot Module Replacement (HMR). Configurat amb un proxy que redirigeix les peticions \texttt{/api} al backend Express, eliminant problemes de CORS en desenvolupament.
	
	\item \textbf{Tailwind CSS v4}: Framework CSS utility-first que genera classes d'utilitat a partir d'una única directiva d'importació. 
	
	\item \textbf{React Router v7}: Llibreria de routing client-side que permet definir rutes protegides amb guards de navegació i sincronitzar filtres amb la URL mitjançant \texttt{useSearchParams}.
	
	\item \textbf{Axios}: Client HTTP amb interceptors que automatitzen la injecció del token JWT a cada petició i la gestió centralitzada d'errors 401 (sessió expirada).
	
	\item \textbf{Recharts}: Llibreria de visualització de dades construïda sobre components React i D3.js. Proporciona gràfics declaratius (donut, àrees apilades, barres, dispersió) amb responsivitat integrada.
	
	\item \textbf{Leaflet + react-leaflet}: Motor de mapes interactius de codi obert amb un pes de ~42 KB. S'han integrat els plugins \texttt{leaflet.markercluster} (agrupació visual) i \texttt{leaflet.heat} (mapes de calor).
	
\end{itemize}

\subsection{Base de dades i infraestructura}

\begin{itemize}
	\item \textbf{PostgreSQL}: Sistema de gestió de bases de dades relacional. Ofereix suport nadiu per a UUID, enums, transaccions ACID i funcions avançades com \texttt{date\_trunc} per a agregacions temporals.
	
	\item \textbf{PostGIS}: Extensió geoespacial per a PostgreSQL que permet emmagatzemar punts geogràfics com a objectes binaris optimitzats (\texttt{geography(Point, 4326)}), que també s'utilitzen en GPS i OpenStreetMap.
	
	\item \textbf{Supabase}: Plataforma Backend-as-a-Service que proporciona una instància de PostgreSQL gestionada amb extensions activables (PostGIS). \cite{9}
	
\end{itemize}

\subsection{Control de versions i gestió del projecte}

\begin{itemize}
	\item \textbf{GitHub}: Plataforma d'allotjament del repositori remot. Facilita la traçabilitat dels canvis, la recuperació de codi i la col·laboració.
\end{itemize}

\section{Desenvolupament}

\subsection{Estructura backend}

L'arquitectura del backend segueix el patró Routes - Middleware - Controller -
Service - Database, on cada capa té una responsabilitat clarament definida.
Aquesta separació de responsabilitats facilita la mantenibilitat i l'escalabilitat
del sistema, ja que permet modificar la lògica de negoci sense afectar les
capes de presentació o d'accés a dades.

\begin{figure}[!h]
	\centering
	\includegraphics[width=0.4\textwidth]{img/architecture}
	\caption{Arquitectura Fontend web/Backend}
\end{figure}

\subsubsection{Arquitectura de la Base de Dades}

La base de dades és PostgreSQL allotjada a Supabase, amb l'extensió PostGIS
per a les capacitats geoespacials. L'esquema es defineix amb Prisma ORM, que
proporciona tipatge automàtic, migracions versionades i un client generat.

S'ha definit un esquema robust amb les seguents taules principals:

\begin{itemize}
	\item \textbf{users}: Emmagatzema els usuaris del sistema amb rol (STUDENT, TECHNICAL,
	ADMIN), punts de gamificació i estat actiu. Inclou el vincle optional amb
	la invitació que va originar l'accés per als rols privilegiats.
	\item \textbf{reports}: Les incidències reportades, amb estat governat per la màquina
	XState, prioritat, categoria tipificada, coordenades GPS, vincles amb
	el creador i el tècnic assignat, i marca temporal de resolució.
	\item \textbf{invites}: Les invitacions per a registre privilegiat, amb token criptogràfic, estat (Pendent/Utilitzada/Revocada) i relació amb l'usuari creat.
\end{itemize}


\subsubsection{Maquina d'Estats}

Per governar el cicle de vida de les incidències (RF-04), s'ha implementat
una màquina d'estats finits utilitzant la llibreria XState v5 amb l'API
setup(). Aquesta decisió respon a la necessitat de centralitzar i formalitzar
les regles de transició, evitant condicions disperses pel codi que serien
difícils de mantenir i verificar.

La màquina defineix cinc estats (Oberta, Assignada, En procés, Validada i
Tancada) i sis transicions, cadascuna protegida per guards que verifiquen
el rol de l'usuari que sol·licita el canvi. Per exemple, només un administrador
pot tancar definitivament una incidència, mentre que assignar-la pot fer-ho
qualsevol usuari que no sigui estudiant.

La integració amb la base de dades és especialment rellevant: quan un servei
necessita executar una transició, crea un actor XState, restaurant l'estat
actual des de la base de dades amb resolveState(). Abans d'executar la
transició, valida amb snapshot.can() si és permesa, permetent retornar
missatges d'error descriptius sense haver de modificar les dades.

\begin{figure}[H]
	\centering
	\includegraphics[width=0.9\columnwidth]{img/transition}
	\caption{Executar transició}
\end{figure}
\subsubsection{Middleware: Xifratge i RBAC}

Per controlar la seguretat del backend i restringir l'acces a determinats usuaris s'ha afegit un middleware que actua com a capa intermitja, filtrant totes aquelles peticions malicioses o invalides i controlant l'acces dels usuaris. 

La primera part del middleware (authenticate) s'encarrega de verificar l'inici de sessió d'un usuari lícit. En rebre una petició, extreu el token del header Authorization, verifica la seva signatura criptogràfica mitjançant jwt.verify(). Si es valid vol dir que l'usuari es correcte i es retorna el payload descodificat (identificador d'usuari i rol). En cas contrari es rebutja amb un codi 401.

La segona part del middleware (authorize) implementa el control d'accés basat en rols
(RBAC). Verifica si l'usuari autenticat posseeix un dels rols requerits. En cas contrari, retorna un codi 403. 

Per al registre, la contrasenya es protegeix amb bcrypt. Per a
l'inici de sessió, es compara la contrasenya proporcionada amb el hash
emmagatzemat i, si coincideix, es genera un token JWT amb una expiració de
24 hores. En ambdós casos, la contrasenya mai es retorna al client, complint
amb els principis de protecció de dades del RGPD.

\begin{figure}[H]
	\centering
	\includegraphics[width=0.9\columnwidth]{img/middleware}
	\caption{Flux de peticions HTTP}
\end{figure}

\subsubsection{Seguretat}

S'ha implementat el patró de registre per invitació, inspirat en el model
Zero Trust: el registre d'estudiants és públic però el backend ignora el
camp role i força STUDENT programàticament. Per als rols privilegiats (ADMIN
i TECHNICAL), es requereix un token d'invitació generat prèviament per un
administrador.

La creació de l'usuari privilegiat i el marcatge de la invitació com a
utilitzada s'executen dins d'una transacció atòmica de Prisma. Si una de les
dues operacions falla, cap de les dues s'aplica, prevenint inconsistències
com un token reutilitzable o un token perdut.

Per prevenir situacions irreversibles, s'han establert dues regles de
protecció: 

\begin{itemize}
	\item La primera protegeix un administrador root (configurat via variable d'entorn)
	que no pot ser revocat sota cap circumstància, garantint que sempre existeixi
	almenys un punt d'accés al sistema.
	
	\item La segona impedeix revocar l'últim administrador actiu mitjançant un comptatge
	previ a la base de dades. Si només resta un administrador actiu, l'operació
	es cancel·la, protegint el sistema de possibles deadlocks.
	
\end{itemize}

\subsection{Panell de control}

L'aplicació consta de cinc pàgines principals, organitzades dins d'un layout
comú amb una barra lateral de navegació fixa i una zona de contingut principal.

\subsubsection{Dashboard}
S'ha escollit Recharts com a llibreria de visualització per la seva integració
nativa amb React. Les diferents gràfiques implementades s'han dividit en diferents blocs: 

\begin{itemize}
	\item El Bloc Operatiu respon la pregunta «Quin és l'estat actual del sistema?»
	mitjançant targetes KPI clicables i un gràfic de rosca que mostra la
	distribució proporcional dels cinc estats.
	\item El Bloc Temporal analitza les tendències amb un gràfic d'àrees apilades que
	mostra simultàniament l'evolució global i la composició per categoria, i un
	gràfic de barres agrupades que compara les incidències creades amb les
	tancades per setmana.
	\item El Bloc de Rendiment avalua l'eficiència dels recursos amb un gràfic de
	barres horitzontals apilades per tècnic i un gràfic de dispersió que relaciona el temps de resolució amb la prioritat.
	\item El Bloc Social mostra un rànquing dels 10 usuaris que més
	incidències han reportat.
\end{itemize}

\subsubsection{Administració d'usuaris}
S'ha implementat una pagina incial d'inici de sessió que permet només als usuaria administradors accedir al panell de control. L'accés es verifica tant al frontend (comporvació rol ADMIN) i al backend mitjançant el middleware (moduls authenticate i authorize).


\subsubsection{Visualització geografica i SIG}
S'ha habilitat l'extensió PostGIS a Supabase per afegir capacitats
geoespacials a la base de dades PostgreSQL. S'ha escollit Leaflet com a motor de mapes i s'ha creat un trigger PL/pgSQL que transforma automàticament els camps decimals latitude i
longitude en un punt binari PostGIS. 

S'han integrat dos plugins específics: leaflet.markercluster per a l'agrupació
visual de markers per proximitat amb animacions de dispersió, i leaflet.heat
per a la generació de mapes de calor amb gradient configurable. 

Aquests plugins en permeten visualitzar en el mapa de manera ordenada els punts on s'ubiquen les incidencies i s'agrupen per proximitat (clustering). 

D'altra banda també s'ha implementat un mapa de calor. Es pot seleccionar el criteri de pes entre tres opcions (prioritat, densitat i antiguitat) per obtenir perspectives analítiques diferents. 

\begin{figure}[H] % Recorda \usepackage{float} al preàmbul
	\centering
	% Fem que ocupi el 100% de la columna, no de la pàgina
	\includegraphics[width=\columnwidth]{img/heatmap} 
	\caption{Heatmap}
\end{figure}

\subsection{Frontend app}

\subsubsection{Inici de sessió i registre}

